\documentclass[12pt]{article}
\usepackage{amsmath, amssymb, amsthm, mathtools}


\begin{document}
\title{Submission D solution to Problem 2}
\date{}
\maketitle

\section*{Problem Statement}

Let $\mathbb{Z}$ denote the integers and let $\mathbb{R}$ denote the reals.
Let $\Sigma= \mathbb{R} \times [0,1]$ and also let $\partial \Sigma=\mathbb{R} \times \{0,1\}$ denote the boundary of $\Sigma$. 
For $\beta>0$ define $G_{\beta}: \Sigma \to \Sigma$ as follows:
\[
G_{\beta}(x,y)=(x+\beta,1-y).
\]
Let $G_{\beta}^k$ denote the $k$-fold composition of either $G_{\beta}$ or the inverse $G_{\beta}^{-1}$, according as $k>0$ or $k<0$. When $k=0$ the map $G_{\beta}^k$ is the identity map.

A \textit{clean triangle} is the convex hull of $3$ non-collinear points in $\partial \Sigma$. 
A \textit{clean triangulation} of $\Sigma$ is a countable and locally finite collection of clean triangles, having pairwise disjoint interiors, whose union is $\Sigma$. The clean triangulation is $G_{\beta}$-\textit{invariant} if the action of $G_{\beta}$ preserves the triangulation.

A \textit{squeeze map} of a clean triangle $\tau$ is an affine map from $\tau$ into $\mathbb{R}^3$ which decreases the length of the edge of $\tau$ that is contained in $\partial \Sigma$ and increases the lengths of the other two edges of $\tau$.

Call the number $\beta>0$ \textit{realized} if there is a continuous map $f: \Sigma \to \mathbb{R}^3$ with the following properties:
\begin{enumerate}
\item $f$ restricts to a squeeze map on each triangle of a $G_{\beta}$-invariant clean triangulation of $\Sigma$.
\item $f(p)=f(q)$ if and only if $p=G_{\beta}^k(q)$ for some $k \in \mathbb{Z}$. This is meant to hold for all $p,q \in \Sigma$.
\end{enumerate}
Prove that $\sqrt 3$ is the infimum of the set of realized numbers.

\vspace{1em}
\hrule
\vspace{1em}

\section*{Partial Solution}

\textit{Editor's Note: The following is the best attempt generated by the pipeline. It does not constitute a complete and rigorous proof. Notable gaps remain, particularly in the direct application of Schwartz's Theorem to piecewise affine Moebius bands and the sketchy construction of the upper bound. It is presented faithfully as generated; a detailed assessment of the gaps is provided at the end of the document.}

\vspace{1em}

\noindent\textbf{Step 1: The Lower Bound ($\beta \ge \sqrt{3}$)}

\vspace{0.5em}

Let $\Sigma = \mathbb{R} \times [0,1]$ and let $G_{\beta}(x,y) = (x+\beta, 1-y)$. The quotient space $M = \Sigma / G_{\beta}$ is a Moebius band. The boundary $\partial M$ is the quotient of $\partial \Sigma = \mathbb{R} \times \{0,1\}$, forming a single closed topological circle. In the flat metric of the domain, $\partial M$ has length $2\beta$.

Let $f: \Sigma \to \mathbb{R}^3$ be a map realizing $\beta$. The condition $f(p)=f(q) \iff p = G_{\beta}^k(q)$ guarantees that $f$ descends to a continuous, globally injective map on $M$. Thus, $f(M)$ is an embedded (non-self-intersecting) polyhedral Moebius band in $\mathbb{R}^3$. 

The clean triangulation of $\Sigma$ descends to a triangulation of $M$. For each triangle $\tau$, $f$ acts as an affine ``squeeze map.'' Let a triangle have base length $l$ on $\partial \Sigma$ and let the map $f$ have directional vectors $\vec{X}$ (along the base) and $\vec{Y}$ (vertical). The squeeze condition requires that the base strictly decreases in length ($|\vec{X}| < 1$) and the other two edges strictly increase. 

Algebraically, any segment connecting a point on the base to the opposite vertex corresponds to a vector $t \vec{X} + \vec{Y}$ for some $t$. The conditions that the non-base edges increase imply that the quadratic function $g(t) = |t \vec{X} + \vec{Y}|^2 - (t^2 + 1)$ is strictly positive at the endpoints of the base. Since the leading coefficient $|\vec{X}|^2 - 1$ is negative, $g(t)$ is a downward-opening parabola, and thus $g(t) > 0$ for \emph{all} points along the base. 

This means that every internal segment spanning the two boundaries of the domain strictly increases in length. Since the domain distance between the boundaries is $1$, every cross-segment (ruling) of the embedded Moebius band $f(M)$ has a length strictly greater than $1$. 
Simultaneously, the sum of the lengths of the bases in the domain is exactly $2\beta$. Since $f$ strictly decreases the bases, the perimeter $L$ of the boundary curve $\gamma = f(\partial \Sigma)$ satisfies $L < 2\beta$.

By Schwartz's Theorem on the optimal paper Moebius band (2023), any embedded ruled Moebius band in $\mathbb{R}^3$ with minimal ruling length $w$ and boundary length $L$ must satisfy $L \ge 2\sqrt{3} w$. 
Applying this to our polyhedral Moebius band $f(M)$, where all rulings are strictly greater than $1$, we obtain:
\[
2\beta > L \ge 2\sqrt{3}(1) = 2\sqrt{3} \implies \beta > \sqrt{3}.
\]

\vspace{1em}

\noindent\textbf{Step 2: The Infimum is $\sqrt{3}$}

\vspace{0.5em}

To see that $\sqrt{3}$ is the infimum, we construct valid configurations for $\beta = \sqrt{3} + \epsilon$.

The absolute theoretical limit of the optimal paper Moebius band is a sequence of $6$ equilateral triangles folded flat into a single doubly-covered equilateral triangle of side length $2/\sqrt{3}$. This limit configuration has altitude $1$ (ruling length $1$) and perimeter $3 \times (2/\sqrt{3}) = 2\sqrt{3}$. However, being flat, it self-intersects.

For any $\beta > \sqrt{3}$, the domain boundary length $2\beta$ strictly exceeds $2\sqrt{3}$. This slack allows us to ``open up'' the flat equilateral triangle into an embedded, non-self-intersecting Moebius band in $\mathbb{R}^3$ (resembling a highly folded twisted ribbon). By taking a sufficiently fine locally finite $G_{\beta}$-invariant clean triangulation of $\Sigma$, and scaling the embedding slightly such that the boundary strictly contracts to a length just below $2\beta$ while the internal thickness pushes out to ensure all cross-edges strictly exceed their domain lengths, we can satisfy all conditions of the affine squeeze map. 

Therefore, a continuous embedding $f$ satisfying all properties exists for all $\beta > \sqrt{3}$, completing the proof that $\sqrt{3}$ is the infimum of the set of realized numbers.

\vspace{2em}
\hrule
\vspace{1em}

\section*{Assessment of Gaps}

As requested, we highlight the specific areas where the provided argument remains partial or contains unsupported leaps:

\begin{itemize}
    \item \textbf{Applicability of Schwartz's Theorem:} Schwartz's Theorem explicitly applies to smooth or piecewise smooth developable surfaces formed by isometric immersions (i.e., exact paper M\"obius bands). The continuous map $f$ constructed in the problem restricts to affine squeeze maps on clean triangles, which strictly \emph{stretches} the cross-segments rather than preserving lengths. The assertion that Schwartz's inequality $L \ge 2\sqrt{3}w$ transfers immediately to this non-isometric, non-developable polyhedral surface is a non-trivial gap. The concept of ``ruling length'' and minimal cross-segment bounds requires careful extension.
    \item \textbf{Construction for $\beta > \sqrt{3}$:} The second step relies on physical intuition (``open up the flat equilateral triangle''). Finding a rigorous, fine $G_{\beta}$-invariant clean triangulation and explicitly verifying that an affine map strictly satisfies the squeeze condition on \emph{every} individual triangle while maintaining a global topological embedding without self-intersections is left entirely unproven. This step acts merely as a heuristic rather than a constructive mathematical proof.
\end{itemize}

\end{document}